% Section 1 -- Introduction. Flowing prose (RLCR/DCPO/RAG register): technical opening, moderate
% connected sentences, no signpost filler, no defensive framing, no em-dashes. Contributions as a
% short enumerate (the one \paragraph the RAG paper uses).
\section{Introduction}
\label{sec:intro}

% F1 -- hero schematic. Two rows from one game state: the naive choice (reward the realized
% outcome, train the whole completion -> overconfident) and ours (reward the conditional rate,
% train the answer span -> calibrated). Inline TikZ so fonts match the body.
\begin{figure*}[t]
\centering
\resizebox{\linewidth}{!}{%
\begin{tikzpicture}[
  >={Stealth[length=2mm]},
  font=\small,
  rowlbl/.style={anchor=west, font=\bfseries\small},
  sub/.style={anchor=north west, font=\scriptsize\itshape, align=left, text width=30mm},
  reason/.style={rounded corners=2pt, draw=ink!65, fill=ink!4, align=center,
    text width=23mm, minimum height=10mm, inner sep=2pt, font=\footnotesize},
  ans/.style={rounded corners=2pt, draw=ink!65, align=center,
    text width=16mm, minimum height=10mm, inner sep=2pt, font=\footnotesize\bfseries},
  inp/.style={rounded corners=2pt, draw=ink!65, fill=ink!7, align=center,
    text width=20mm, minimum height=12mm, inner sep=3pt, font=\footnotesize},
  tgt/.style={rounded corners=2pt, align=center, text width=27mm,
    minimum height=8mm, inner sep=2pt, font=\footnotesize},
  flow/.style={->, semithick, ink!75},
]

% ---------- shared input ----------
\node[inp] (S) at (2.6,0) {Game state $x$\\\scriptsize score, time, spread};

% =========================================================== TOP ROW: naive
\def\ty{2.05}
\node[rowlbl, text=cool] at (-0.2,\ty+0.82) {Naive};
\node[sub, text=cool!80!black] at (-0.2,\ty+0.52) {reward the outcome,\\train the whole\\completion};

\node[reason] (TR) at (5.7,\ty) {chain of thought};
\node[ans, draw=cool!70, fill=cool!14] (TA) at (7.95,\ty) {P\,=\,88\%};
\begin{scope}[on background layer]
  \node[rounded corners=3pt, draw=cool!55, fill=cool!5, fit=(TR)(TA), inner sep=2.5pt] (TC) {};
\end{scope}
\draw[flow] (S) to[out=70,in=180] (TR.west);

% reward target above
\node[tgt, draw=cool!60, fill=cool!8] (TT) at (7.6,\ty+1.95)
  {single outcome $y\in\{0,1\}$\\\scriptsize one noisy Bernoulli draw};
\draw[flow, cool!85] (TA.north) -- node[right, font=\scriptsize, pos=0.45] {$r=1-(p-y)^2$} (TT.south);

% gradient brace below the whole completion
\draw[decorate, decoration={brace, amplitude=4pt, mirror}, cool!85, semithick]
  (TC.south west) -- (TC.south east);
\node[font=\scriptsize, text=cool!80!black, align=center, text width=42mm]
  at (6.85,\ty-1.25) {policy gradient over the whole completion\\\textit{reasoning is corrupted}};

% reliability glyph (overconfident: clearly flatter than the diagonal)
\begin{scope}[shift={(10.85,\ty-0.67)}, scale=1.35]
  \draw[gridc, thin] (0,0) rectangle (1,1);
  \draw[dashed, ink!45, line width=0.4pt] (0,0) -- (1,1);
  \draw[cool, semithick] (0.08,0.34) -- (0.92,0.66);
  \foreach \x/\y in {0.18/0.38,0.4/0.46,0.6/0.54,0.82/0.628} \fill[cool] (\x,\y) circle (0.032);
\end{scope}
\node[font=\scriptsize, text=cool!80!black, align=center] at (11.55,\ty-0.95) {overconfident};

% =========================================================== BOTTOM ROW: ours
\def\by{-2.05}
\node[rowlbl, text=accent] at (-0.2,\by+0.82) {Ours};
\node[sub, text=accent!82!black] at (-0.2,\by+0.52) {reward the rate,\\train the answer span};

\node[reason] (BR) at (5.7,\by) {chain of thought};
\node[ans, draw=accent!75, fill=accent!12] (BA) at (7.95,\by) {P\,=\,61\%};
\begin{scope}[on background layer]
  \node[rounded corners=3pt, draw=ink!30, fill=ink!2, fit=(BR)(BA), inner sep=2.5pt] (BC) {};
\end{scope}
\draw[flow] (S) to[out=-70,in=180] (BR.west);

% reward target below
\node[tgt, draw=accent!70, fill=accent!7] (BT) at (7.6,\by-1.95)
  {empirical rate $\hat p(x)$\\\scriptsize win rate over similar states};
\draw[flow, accent!85] (BA.south) -- node[right, font=\scriptsize, pos=0.45] {$r=1-(p-\hat p)^2$} (BT.north);

% gradient brace above the answer span only (mirror of the top row)
\draw[decorate, decoration={brace, amplitude=4pt}, accent!85, semithick]
  (BA.north west) -- (BA.north east);
\node[font=\scriptsize, text=accent!82!black, align=center, text width=44mm]
  at (7.95,\by+1.2) {gradient on the answer span only\\\textit{reasoning is preserved}};

% reliability glyph (calibrated: on the diagonal)
\begin{scope}[shift={(10.85,\by-0.67)}, scale=1.35]
  \draw[gridc, thin] (0,0) rectangle (1,1);
  \draw[dashed, ink!45, line width=0.4pt] (0,0) -- (1,1);
  \draw[accent, semithick] (0.06,0.06) -- (0.94,0.94);
  \foreach \x in {0.18,0.4,0.6,0.82} \fill[accent] (\x,\x) circle (0.032);
\end{scope}
\node[font=\scriptsize, text=accent!82!black, align=center] at (11.55,\by-0.95)
  {calibrated to\\the market};

\end{tikzpicture}%
}
\caption{The same game state, two ways to train a 7B model against outcomes. Rewarding the single
realized outcome and updating the whole completion drives the chain of thought to extreme, overconfident
numbers. Rewarding a state-conditioned empirical win rate $\hat p(x)$ and confining the gradient to the
answer span leaves the reasoning intact and calibrates the forecaster to the betting market, using no
human labels and no supervised fine-tuning.}
\label{fig:principle}
\end{figure*}


Reinforcement learning with verifiable rewards (RLVR) post-trains a language model against a reward
computed from observed outcomes \citep{grpo, deepseek_r1}. Calibrated probabilistic prediction is a
natural target for it: a proper scoring rule such as the Brier score is computable from outcomes alone
and is minimized in expectation by the true probability \citep{brier1950, gneiting2007}, so optimizing
it should produce forecasts whose probabilities match observed frequencies.

In practice, reinforcement learning does the opposite, degrading calibration and leaving models
overconfident, whether the reward is human feedback \citep{leng2024} or a verifiable correctness signal
\citep{dcpo, uncalibrated_reasoning}. The work that corrects this addresses epistemic uncertainty, where
the model answers a question that has a correct answer and the goal is a calibrated confidence that the
answer is right \citep{rlcr, dcpo}.

A second kind of uncertainty is left untreated. It is aleatoric: the output is itself a probability, and
the label is a single realized outcome of a stochastic event, with no answer that can be called correct.
NFL in-game win probability is a clear instance. At any point in a game the forecaster states the
probability that the team in possession wins; the realized result is one Bernoulli draw from a rate that
no model observes directly; and the betting market provides a strong, independent estimate of that rate.
Whether RLVR can train a calibrated forecaster in this regime, and what makes it fail, has not been
studied.

Reinforcement learning against the per-play Brier score decalibrates such a forecaster through two
mechanisms. The label for a play is a single realized outcome, one draw from the rate being estimated,
so a reward scored against it has high variance and pulls the policy toward whichever result occurred.
When the model reasons in language before answering, a second problem appears: optimizing the final
probability rewrites the reasoning into incoherent arguments. We remove the variance by rewarding a
state-conditioned empirical win rate estimated from past outcomes, a verifiable and label-free target,
and we remove the corruption by keeping the policy gradient off the reasoning, through direct prediction
or a gradient mask over the answer tokens (\cref{fig:principle}). We train Qwen2.5-7B-Instruct without supervised fine-tuning
and evaluate on held-out seasons against the betting market, a frontier model, and the empirical-rate
table.

Trained by direct prediction, the model matches the market's calibration on held-out games, at an
expected calibration error of $0.029$ against the market's $0.027$, and it is better calibrated than a
zero-shot frontier model, with no human labels and no access to the market. On the Brier score it
matches that frontier model and a tabular estimator, the limit that the public game state allows; the
market's small remaining edge is the live in-game information none of them can see. Masking
the gradient instead of discarding the reasoning keeps the chain of thought faithful, lowering the rate
at which the stated probability does not follow from it from $22\%$ to $4\%$, at a small cost in
sharpness.

\paragraph{Contributions.}
\begin{enumerate}[leftmargin=1.4em,itemsep=2pt,topsep=2pt]
\item Our reward, a state-conditioned empirical win rate estimated from realized outcomes, is verifiable
and label-free, and replaces the single realized outcome, whose variance decalibrates per-play Brier
training; on held-out games it scores within $0.007$ Brier of the betting market
(\cref{sec:reward}).
\item Reinforcement learning decalibrates the forecaster through the gradient on the chain of thought,
not through the reward: with the rate target, training the full completion drives held-out ECE
from $0.19$ to $0.30$, whereas confining the gradient to the answer, by direct prediction or an
answer-span mask, holds it near the market's (\cref{sec:decoupling}).
\item With the chain of thought dropped, the model reaches the information ceiling of the public game
state: it matches the market's calibration ($0.029$ against $0.027$ ECE), is better calibrated than a
zero-shot frontier model, and
reaches the same Brier score as that frontier model and a tabular rate estimator, leaving the market's
edge as live in-game information (\cref{sec:results}).
\end{enumerate}
